\section{Status, authorship, and separation from the reader}

The accompanying critical reader asks what a named source proves, what it
does not prove, and which repairs are required to state its mathematics
correctly.  This document asks a different question: what exact mathematics
can be proved after the source conventions and implementation coordinates have
been reconstructed?  Results proved here are therefore not silently
attributed to Barina, Oliveira e Silva, Roosendaal, Hercher, or the authors of
the public code.

The status vocabulary is literal.

\begin{description}[leftmargin=3.4cm,style=nextline]
\item[Proved.] The displayed hypotheses imply the conclusion by the proof
printed here.
\item[Finite certificate.] A proved reduction leaves a finite set, and the
identified deterministic program exhausts exactly that set.  The bound and
the items not certified are stated explicitly.
\item[Source report.] A source says that a historical computation occurred.
The statement is evidence about the reported run, not a reproduction of it.
\item[Open.] No proof or counterexample has been established in this project.
An open item is never used downstream as though it were a lemma.
\end{description}

Two map conventions recur.  The unshortened and shortened Collatz maps are
\[
 U(n)=\begin{cases}3n+1,&n\text{ odd},\\n/2,&n\text{ even},\end{cases}
 \qquad
 T(n)=\begin{cases}(3n+1)/2,&n\text{ odd},\\n/2,&n\text{ even}.
 \end{cases}
\]
They have different clocks.  One odd $T$-step is two $U$-steps; one even
$T$-step is one $U$-step.  No statement below identifies their iterate
indices.

\begin{nonclaim}
Nothing in this document proves that every positive Collatz orbit reaches
$1$.  A dead sieve bit means that one of the stated finite certificates
applies; a live bit means only that the implemented certificate family did not
eliminate it.
\end{nonclaim}

